% GENERATED FILE. Edit canonical lessons, then run npm run generate:books.
% canonical-source-hash: fnv1a64:0ddb2d073f61d38a
% canonical-lessons: GU-C511-majur, GU-C511-mukadam, GU-C511-kadiyo, GU-C511-mistri, GU-C511-karkun

\chapter{People at Work}
\label{ch:gu-people-at-work}
\hlchaptermodality{511}

\begin{chapteropening}
\textbf{By the end of this chapter:} \emph{I can name the people who work on a building site and in an office.}

Complete the last lesson of chapter 511: I can name the people who work on a building site and in an office.
\end{chapteropening}

\section[\texorpdfstring{\gu{મજૂર} (majūr)}{મજૂર (majūr)}]{\gu{મજૂર} (majūr) — a labourer}

\label{lesson:GU-C511-majur}



\begin{quote}
Before the new one: say the Gujarati for to cast, then the Gujarati for to have (something) made.
\end{quote}

\subsection*{You'll want to know: \gu{મજૂર}}
\textbf{\gu{મજૂર}} — \emph{majūr} — \textquotedblleft{}a labourer\textquotedblright{}.

\begin{quote}
\gu{મજૂરો} \gu{સવારથી} \gu{સાંજ} \gu{સુધી} \gu{કામ} \gu{કરે} \gu{છે}.
\end{quote}

\begin{quote}
\emph{majūro savārthī sānj sudhī kām kare che.}
\end{quote}

\begin{quote}
The labourers work from morning till evening.
\end{quote}

A masculine or feminine noun, by the person meant.

\begin{cousinweb}
\textbf{\gu{મજૂર}} (\emph{majūr}) is usually traced to Persian \emph{mazdūr}, \textquotedblleft{}one who works for a wage\textquotedblright{}.
\end{cousinweb}

\begin{grammarlens}[title={What you've built}]
One more word for describing what you lived and felt.
\end{grammarlens}

\subsection*{Your turn}
\begin{itemize}
\raggedright
  \item \emph{Say it:} \emph{majūr}
  \item \emph{Say it:} \emph{majūr}, once more
  \item \emph{Say it:} \emph{majūro savārthī sānj sudhī kām kare che.} — The labourers work from morning till evening.
  \item \emph{Recall:} say the Gujarati for to cast, then the Gujarati for to have (something) made, then say \emph{majūr} again
\end{itemize}

\begin{tcolorbox}[breakable,colback=teal!4,colframe=teal!35!black,title={Before you move on}]
What does \gu{મજૂર} mean? (\textquotedblleft{}A labourer\textquotedblright{}.) Say it once more.
\end{tcolorbox}
\section[\texorpdfstring{\gu{મુકાદમ} (mukādam)}{મુકાદમ (mukādam)}]{\gu{મુકાદમ} (mukādam) — a foreman}

\label{lesson:GU-C511-mukadam}



\begin{quote}
Before the new one: say the Gujarati for to have (something) made, then the Gujarati for a labourer.
\end{quote}

\subsection*{You'll want to know: \gu{મુકાદમ}}
\textbf{\gu{મુકાદમ}} — \emph{mukādam} — \textquotedblleft{}a foreman\textquotedblright{}.

\begin{quote}
\gu{મુકાદમે} \gu{મજૂરોને} \gu{કામ} \gu{સમજાવ્યું}.
\end{quote}

\begin{quote}
\emph{mukādme majūrone kām samjāvyũ.}
\end{quote}

\begin{quote}
The foreman explained the work to the labourers.
\end{quote}

A masculine noun.

\begin{cousinweb}
\textbf{\gu{મુકાદમ}} (\emph{mukādam}) is usually traced to Arabic \emph{muqaddam}, \textquotedblleft{}one placed in front, a headman\textquotedblright{}, which came to Gujarati through Persian.
\end{cousinweb}

\begin{grammarlens}[title={What you've built}]
One more word for describing what you lived and felt.
\end{grammarlens}

\subsection*{Your turn}
\begin{itemize}
\raggedright
  \item \emph{Say it:} \emph{mukādam}
  \item \emph{Say it:} \emph{mukādam}, once more
  \item \emph{Say it:} \emph{mukādme majūrone kām samjāvyũ.} — The foreman explained the work to the labourers.
  \item \emph{Recall:} say the Gujarati for to have (something) made, then the Gujarati for a labourer, then say \emph{mukādam} again
\end{itemize}

\begin{tcolorbox}[breakable,colback=teal!4,colframe=teal!35!black,title={Before you move on}]
What does \gu{મુકાદમ} mean? (\textquotedblleft{}A foreman\textquotedblright{}.) Say it once more.
\end{tcolorbox}
\section[\texorpdfstring{\gu{કડિયો} (kaḍiyo)}{કડિયો (kaḍiyo)}]{\gu{કડિયો} (kaḍiyo) — a mason}

\label{lesson:GU-C511-kadiyo}



\begin{quote}
Before the new one: say the Gujarati for a labourer, then the Gujarati for a foreman.
\end{quote}

\subsection*{You'll want to know: \gu{કડિયો}}
\textbf{\gu{કડિયો}} — \emph{kaḍiyo} — \textquotedblleft{}a mason\textquotedblright{}.

\begin{quote}
\gu{કડિયાએ} \gu{ઈંટોથી} \gu{નવી} \gu{દીવાલ} \gu{બનાવી}.
\end{quote}

\begin{quote}
\emph{kaḍiyāe ī̃ṭothī navī dīvāl banāvī.}
\end{quote}

\begin{quote}
The mason built a new wall of bricks.
\end{quote}

A masculine noun.

\begin{cousinweb}
\textbf{\gu{કડિયો}} (\emph{kaḍiyo}) builds with brick, stone and mortar. As the doer of a past action it becomes \textbf{\gu{કડિયાએ}}. No dictionary we use records an older form, so its history is not settled.
\end{cousinweb}

\begin{grammarlens}[title={What you've built}]
One more word for describing what you lived and felt.
\end{grammarlens}

\subsection*{Your turn}
\begin{itemize}
\raggedright
  \item \emph{Say it:} \emph{kaḍiyo}
  \item \emph{Say it:} \emph{kaḍiyo}, once more
  \item \emph{Say it:} \emph{kaḍiyāe ī̃ṭothī navī dīvāl banāvī.} — The mason built a new wall of bricks.
  \item \emph{Recall:} say the Gujarati for a labourer, then the Gujarati for a foreman, then say \emph{kaḍiyo} again
\end{itemize}

\begin{tcolorbox}[breakable,colback=teal!4,colframe=teal!35!black,title={Before you move on}]
What does \gu{કડિયો} mean? (\textquotedblleft{}A mason\textquotedblright{}.) Say it once more.
\end{tcolorbox}
\section[\texorpdfstring{\gu{મિસ્ત્રી} (mistrī)}{મિસ્ત્રી (mistrī)}]{\gu{મિસ્ત્રી} (mistrī) — a mechanic}

\label{lesson:GU-C511-mistri}



\begin{quote}
Before the new one: say the Gujarati for a foreman, then the Gujarati for a mason.
\end{quote}

\subsection*{You'll want to know: \gu{મિસ્ત્રી}}
\textbf{\gu{મિસ્ત્રી}} — \emph{mistrī} — \textquotedblleft{}a mechanic\textquotedblright{}.

\begin{quote}
\gu{ગાડી} \gu{બગડી} \gu{ગઈ} \gu{એટલે} \gu{અમે} \gu{મિસ્ત્રીને} \gu{બોલાવ્યો}.
\end{quote}

\begin{quote}
\emph{gāḍī bagḍī gaī eṭle ame mistrīne bolāvyo.}
\end{quote}

\begin{quote}
The car broke down, so we called a mechanic.
\end{quote}

A masculine noun.

\begin{cousinweb}
\textbf{\gu{મિસ્ત્રી}} (\emph{mistrī}) is usually traced to Portuguese \emph{mestre}, \textquotedblleft{}a master of a trade\textquotedblright{}. In Gujarati it names a skilled hand: a mechanic, or a master carpenter or mason.
\end{cousinweb}

\begin{grammarlens}[title={What you've built}]
One more word for describing what you lived and felt.
\end{grammarlens}

\subsection*{Your turn}
\begin{itemize}
\raggedright
  \item \emph{Say it:} \emph{mistrī}
  \item \emph{Say it:} \emph{mistrī}, once more
  \item \emph{Say it:} \emph{gāḍī bagḍī gaī eṭle ame mistrīne bolāvyo.} — The car broke down, so we called a mechanic.
  \item \emph{Recall:} say the Gujarati for a foreman, then the Gujarati for a mason, then say \emph{mistrī} again
\end{itemize}

\begin{tcolorbox}[breakable,colback=teal!4,colframe=teal!35!black,title={Before you move on}]
What does \gu{મિસ્ત્રી} mean? (\textquotedblleft{}A mechanic\textquotedblright{}.) Say it once more.
\end{tcolorbox}
\section[\texorpdfstring{\gu{કારકુન} (kārkun)}{કારકુન (kārkun)}]{\gu{કારકુન} (kārkun) — a clerk}

\label{lesson:GU-C511-karkun}



\begin{quote}
Before the new one: say the Gujarati for a mason, then the Gujarati for a mechanic.
\end{quote}

\subsection*{You'll want to know: \gu{કારકુન}}
\textbf{\gu{કારકુન}} — \emph{kārkun} — \textquotedblleft{}a clerk\textquotedblright{}.

\begin{quote}
\gu{કચેરીમાં} \gu{કારકુન} \gu{આખો} \gu{દિવસ} \gu{કાગળો} \gu{લખે} \gu{છે}.
\end{quote}

\begin{quote}
\emph{kacerīmā̃ kārkun ākho divas kāgḷo lakhe che.}
\end{quote}

\begin{quote}
In the office the clerk writes papers all day.
\end{quote}

A masculine noun.

\begin{cousinweb}
\textbf{\gu{કારકુન}} (\emph{kārkun}) is usually traced to Persian \emph{kārkun}, \textquotedblleft{}one who does the work, an agent\textquotedblright{}, built on \emph{kār}, \textquotedblleft{}work\textquotedblright{}.
\end{cousinweb}

\begin{grammarlens}[title={What you've built}]
One more word for describing what you lived and felt.
\end{grammarlens}

\subsection*{Your turn}
\begin{itemize}
\raggedright
  \item \emph{Say it:} \emph{kārkun}
  \item \emph{Say it:} \emph{kārkun}, once more
  \item \emph{Say it:} \emph{kacerīmā̃ kārkun ākho divas kāgḷo lakhe che.} — In the office the clerk writes papers all day.
  \item \emph{Recall:} say the Gujarati for a mason, then the Gujarati for a mechanic, then say \emph{kārkun} again
\end{itemize}

\begin{tcolorbox}[breakable,colback=teal!4,colframe=teal!35!black,title={Before you move on}]
What does \gu{કારકુન} mean? (\textquotedblleft{}A clerk\textquotedblright{}.) Say it once more.
\end{tcolorbox}
